\section{Conclusões}

A presente etapa da iniciação científica teve como objetivo principal a 
consolidação teórica, metodológica e computacional da aplicação de 
algoritmos quânticos variacionais à simulação de sistemas moleculares, 
utilizando o algoritmo Variational Quantum Eigensolver (VQE) como estudo de 
caso. Os resultados obtidos demonstram que a implementação desenvolvida foi 
capaz de reproduzir com elevada precisão a energia do estado fundamental da 
molécula de hidrogênio (H$_2$), validando o correto funcionamento do fluxo 
híbrido quântico--clássico adotado e confirmando a consistência da 
metodologia empregada.

A extensão do procedimento para o sistema LiH evidenciou limitações 
relevantes associadas à escalabilidade do VQE em sistemas moleculares de 
maior complexidade eletrônica. O aumento do erro variacional e a ausência de 
convergência estável indicam que a expressividade do \textit{ansatz} 
\textit{hardware-efficient} e a estratégia de otimização clássica utilizada 
tornam-se fatores críticos à medida que cresce a dimensionalidade do espaço 
variacional. Dessa forma, os resultados não apenas confirmam a viabilidade 
do método em sistemas simples, mas também reproduzem, em ambiente 
controlado, desafios característicos da era NISQ, especialmente aqueles 
relacionados ao equilíbrio entre profundidade de circuito, número de 
parâmetros e estabilidade do processo de otimização.

Sob a perspectiva metodológica, o projeto resultou na implementação completa 
e reprodutível do \textit{pipeline} de simulação molecular baseado em VQE, 
incluindo construção do Hamiltoniano eletrônico, mapeamento férmion--qubit, 
definição de estados variacionais, execução do ciclo híbrido de otimização e 
validação clássica por meio de diagonalização exata. A organização 
sistemática dos experimentos em repositório público versionado estabelece 
uma base computacional consistente para a continuidade do projeto, 
garantindo transparência científica e extensibilidade das simulações 
realizadas.

Os resultados obtidos nesta fase caracterizam, portanto, uma prova de 
conceito bem-sucedida da aplicação de algoritmos variacionais em simulação 
molecular, estabelecendo as condições necessárias para avanços em direção à 
modelagem de sistemas de maior relevância tecnológica associados ao 
desenvolvimento de materiais para armazenamento de energia.

Como continuidade natural do projeto, destacam-se os seguintes 
direcionamentos futuros:

\begin{itemize}
    \item investigação de \textit{ansätze} quimicamente inspirados, como Unitary Coupled Cluster (UCC) e abordagens adaptativas, visando maior capacidade de representação da correlação eletrônica;
    
    \item avaliação de diferentes estratégias de otimização clássica, incluindo métodos baseados em gradiente e técnicas híbridas mais robustas;
    
    \item aplicação de técnicas de mitigação de ruído e execução das simulações em hardware quântico real disponibilizado pela plataforma IBM Quantum;
    
    \item extensão das simulações para moléculas maiores e sistemas eletrônicos mais representativos de materiais relevantes para tecnologias de armazenamento de energia;
    
    \item integração de técnicas de aprendizado de máquina para aceleração do processo variacional, predição de parâmetros iniciais e análise da paisagem de otimização;
    
    \item realização de estudos comparativos sistemáticos entre simulações quânticas variacionais e métodos clássicos consolidados de química computacional.
\end{itemize}

Dessa forma, os resultados obtidos estabelecem uma infraestrutura 
experimental reutilizável. A etapa atual do projeto cumpre o papel de 
estabelecer os fundamentos científicos e computacionais necessários para 
investigações futuras voltadas à utilização integrada de computação quântica 
e aprendizado de máquina na descoberta e otimização de materiais energéticos 
avançados.
